% =============================================================
%  Conclusion.tex —— §12 结论
%  数据来源：四问素材库 + 真机汇总 + 各章节小结
%  编译方式：xelatex Conclusion.tex
% =============================================================
\documentclass[UTF8,12pt,a4paper]{ctexart}

\usepackage{amsmath}
\usepackage{amssymb}
\usepackage{geometry}
\geometry{margin=2.5cm}
\usepackage{booktabs}
\usepackage{float}
\usepackage{caption}
\usepackage{enumitem}
\usepackage{titlesec}
\titlespacing*{\section}{0pt}{4pt}{2pt}
\titlespacing*{\subsection}{0pt}{4pt}{1pt}
\setlength{\parskip}{1pt}
\linespread{1.05}

\setcounter{section}{11}

\begin{document}

% =============================================================
\section{结论}
% =============================================================

% -------------------------------------------------------------
\subsection{逐题结论}
% -------------------------------------------------------------

\textbf{针对问题一}（$n\!=\!15$ 单车 TSP）：以 one-hot 位置编码构建 225 比特 QUBO（罚系数 $A\!=\!200$），由 Held-Karp 精确 DP、Kaiwu SDK SA 与 CIM 真机三层求解器分别得到 $J\!=\!29/29/31$；Held-Karp 与 SDK SA 经 polish 后均达全局最优 $J\!=\!29$，CIM 真机 $J\!=\!31$（gap $+6.9\%$），最优路径见正文路径图。

\textbf{针对问题二}（$n\!=\!15$ 单车 TSPTW）：保留 225 比特 one-hot 编码，时间窗罚不进入 QUBO 而在解码后由统一 evaluate 函数计算，避免线性化导致的可行率坍塌。Python、Kaiwu SDK SA 与 CIM 真机四方一致得 $J\!=\!84121$（$travel\!=\!31$，$penalty\!=\!84090$，违反 12/15 客户），构成跨算法跨硬件的重要稳定性证据，作为本实例的当前最优可行解。

\textbf{针对问题三}（$n\!=\!50$ 大规模单车 TSPTW）：直接构造 $50^2\!=\!2500$ 比特 QUBO 超出真机比特预算，故采用滚动窗口分解为 3 段，子 QUBO 维度分别为 $289/289/256$，均不超过 550。Python LNS、Kaiwu SDK SA 与 CIM 真机三方一致达 $J\!=\!4{,}941{,}906$（$travel\!=\!56$，$penalty\!=\!4{,}941{,}850$）；CIM seg3 经 polish 后达到 Python/SDK 基线，验证了 ``大规模分解 + 真机 + polish'' 流程在本实例中的可行性；在当前实现和参数设置下 SDK 较 Python 取得约 $26\times$ 加速。

\textbf{针对问题四}（$n\!=\!50$ 多车 VRPTW）：采用启发式分配 $+$ 每车独立子 TSP QUBO（$\le\!64$ 比特）的分层编码，绕开 12500 比特直接编码的不可行性，求得近似最优 $J\!=\!7149$（$K^{\ast}\!=\!7$，$travel\!=\!109$，$penalty\!=\!40$）；CIM 真机 $K\!=\!7$ 得 $J\!=\!7168$（gap $+0.27\%$），Python/SDK/CIM 结果接近。$K\!\in\![5,10]$ 灵敏度显示本文在 $K\!\in\!\{5,6,8,9\}$ 严格优于参考，$K\!=\!7$ 持平；$K\!=\!8$ block-diag 副方案揭示了 ``配额 vs 解码质量'' 的真机部署 trade-off。

% -------------------------------------------------------------
\subsection{综合结论}
% -------------------------------------------------------------

\textbf{(1) 方法链路：} 本文建立 ``精确基线 $\to$ QUBO 编码 $\to$ SDK 验证 $\to$ CIM 真机 $\to$ 大规模分解与回评'' 的统一求解链路，以 \texttt{solve\_subqubo(Q, backend)} 抽象层把 SDK 与真机求解收敛为一行后端切换，满足题目对 SDK 验证与真机求解的要求。\textbf{(2) 比特预算：} 四问所有子 QUBO 维度均低于 550（Q1/Q2 为 225、Q3 每段 $\le 289$、Q4 每车 $\le 64$），并经 8-bit 整数精度调整后上传，均满足本文真机提交的规模约束。\textbf{(3) 真实性与可复现：} 全文累计 12 次 CIM 真机配额提交，其中 Q3 seg1 因配额异常未获得有效返回并如实标注；所有真机结果均以平台日志、原始返回与统一 evaluate 回评结果为依据。\textbf{(4) 工程价值：} 滚动窗口分解、每车独立子 QUBO 与 block-diag 合并三类策略形成一套与 550 比特真机上限相协调的求解流程，为类似规模受限的 QUBO 路径规划问题提供了可参考流程。

% -------------------------------------------------------------
\subsection{展望}
% -------------------------------------------------------------

本文工作可在三方向延伸：(1) 引入约束保持邻域或块翻转策略，提升 one-hot QUBO 的原始可行率，降低对解码后 polish 的依赖；(2) 针对 Q2 一类小规模 TSPTW 引入更强的精确算法或 ILP / CP-SAT 验证，以更严格判断 $J\!=\!84121$ 的最优性；(3) 探索客户分配与车内排序的联合 QUBO 编码，分析比特数与求解质量之间的权衡。

\end{document}
